\documentclass[12pt]{article}
\usepackage{styles/cwilliams-standard}

% Set up document metadata
\setauthors{Tzu-Chen Chiu, Arjun Sivakumar, Mara Schimmel, Christopher Williams}
\setheaderauthors{Chiu, Sivakumar, Schimmel, Williams}
\settitle{Final Product Submission}
\setclass{IPD}

\begin{document}

\maketitlepage

\section{Product Description}

\textbf{Background}: Google has published work that details their use of Google's DeepMind machine learning knowledge and Google Quantum AI's error correction expertise to create Google's AlphaQubit.

\textbf{Description}: Reproduce Google's AlphaQubit and create a Recurrent Transformer-based Neural Network for Quantum Error Correction. 

\section{Design}

Our core model is aiming to replicate AlphaQubit's architecture. It is designed to learn temporal and spatial correlations in syndrome data across multiple rounds of error corretion.

\textbf{Model Iputs and Embeddings}: For each shot, we treat the sequence of stabilizer measurements over time as a sequence-like input:

\begin{itemize}
    \item \textbf{Syndrome}: Values (hard 0/1 or soft analog) are first project into a d\_model dimensional embedding using a linear layer.
    \item \textbf{Stab\_id} and \textbf{Cycle\_id} are converted into learnable embeddings:
    \begin{itemize}
        \item \textbf{stab\_id} encodes \textit{which stabilizer} generated the syndrome
        \item \textbf{cycle\_id} encodes \textit{which round} of the error-correction cycle produced that measurement
    \end{itemize}
\end{itemize}

These embeddings are combined (e.g., summed or concatenated + projected) to form a per-site feature vector of dimension d\_model. The full sequence is then processed by transformer encoder layers.

\subsection*{Transformer Architecture}

\begin{itemize}
    \item We use \textbf{binary cross-entropy with logits (BCEWithLogitsLoss)} to train the model to predict the logical error label for each shot.
    \item Optimization is performed using \textbf{AdamW} with:
    \begin{itemize}
        \item Learning rate $1e-3$
        \item Weight decay $1e-5$
    \end{itemize}
\end{itemize}

These hyperparameters are standard for transformer training and were validated by monitoring training/validation loss and the Logical Error Rate (LER) across epochs.

\subsection*{Training, Validation, LER}

We split the dataset into training and validation sets (e.g., 80\% / 20\%) and use PyTorch's DataLoader for mini-batch training.

\begin{itemize}
    \item \textbf{During training}: The model minimizes BCE loss on the training set.
    \item \textbf{During validation}: We evaluate both validation loss and validation LER.
\end{itemize}

LER is computed by:

\begin{enumerate}
    \item Applying a sigmoid to the model's logits to get probabilites.
    \item Thresholding at 0.5 to get binary predictions.
    \item Comparing predictions with the true labels and averaging the mismatch rate.
\end{enumerate}

We save the model-checkpoint with the lowest validation LER to ensure that the selected model is optimized for decoding performance, not just classification loss.

In our current experiments on a distance-3 surface code with soft measurement data, the learned decoder achieves an LER of about 3\%, 
while the classical baseline MWPM achieves roughly 0.3\%, which serves as a reference point and validates our pipeline.

\section{Ethical Implications}

Due to the fact that we are potentially aiding in the development of cryptographic breaking through quantum computers, there may be an impact on the classical 
cryptography scene using the power of quantum computers. This may impact consumer privacy by allowing vastly more efficient decryption of encrypted data, 
as well as governments or financial institutions. 

We are not generating a specific product, rather the technology for decryption. 
Therefore, it is difficult to effectively minimize the negative consequences to different groups, as we do not have much say in its potential 
future uses, nor its reach after it enters the public. 

In terms of biases, there are a few to account for, which will be handled to the best of our abilities. We will use publicly available data, 
or generate our own quantum codes using publicly available simulators. Another potential bias to account for is the model we are basing our 
project on which is used for the surface code, while our project aims to work with codes other than the surface code. This means that we 
may not be able to achieve our secondary goal, which is to adapt the program for other codes.

\section{Current MVP and Stretch Goals}

The MVP for this semester was to begin training the Transformer model, as well as implement a small Minimum Weight Perfect Matching (MWPM) 
code for comparison and data. The stretch goal for this semester was to fully train the model to be comparable to Google’s model. 
We successfully began training the transformer model, in addition to generating a small surface code. We also managed to generate graphs 
comparing physical error rate vs logical error rate and logical error rate for each distance (3, 5, 7, 9), for both the MWPM and Transformer 
model. Although we did hit or MVP for this semester of constructing a transformer model and implementing the MWPM baseline algorithm, we 
were not able to hit out stretch goals as our evaluation and results of the transformer model that we built indicate it does not perform 
as well as Google’s AlphaQubit.

\section{Product Functionalities}

As our product does not have a direct UI/UX, we describe its functionality at a system level. Please refer to our README.md in our 
project repository for how to run our model and reproduce our results.

Our product is a neural quantum error decoder that analyzes stabilizer measurement data
from quantum error correction circuits and predicts the most likely logical errors.
By providing highly accurate decoding decisions, the system enables quantum computers
to identify and mitigate logical failures during computation.

The decoding process runs transparently in the background as part of the quantum
error correction pipeline, without direct user interaction. From the user's
perspective, the system improves computational reliability while remaining invisible.

To support evaluation and analysis, we also provide visualization diagrams that display
the logical error rate (LER) based on user-provided inputs, allowing users to better
understand decoder performance under different error conditions.

\begin{figure}[H]
    \centering
    \includegraphics[width=0.5\textwidth]{images/google_ler.png}
    \caption{Google LER Results at Different Code Distances}
\end{figure}

\section{Semester 2 Sprint 1 Backlog}

For this semester, one of the things that went well was gaining a fundamental knowledge and basic understanding of the quantum computing, 
and quantum error detection solutions.This was a really important achievement as the material is quite dense and new for many of us, 
and it was important to have that fundamental knowledge before going ahead and implementing a neural network error correction algorithm. 
Other things that went well include our consistent team work and regular meetings with our mentor from Northrop Grumman, which really 
helped keep us on track for the project and break down the larger goal into smaller, simpler and achievable tasks. Ultimately, we 
also achieved our MVP goal of building a transformer decoder for this semester. Things that can be improved are focusing more on 
learning about neural network architectures, and how to best build a recurrent transformer network that can achieve comparable 
model performance to Google’s Alpha Qubit. We also need to improve our model pipeline, which involves streamlining simulation and 
evaluation processes, and improving code quality and structure. 

\begin{figure}[H]
    \centering
    \includegraphics[width=0.90\textwidth]{images/sprint_4_backlog.png}
    \caption{Semester 2 Sprint 1 Backlog}
    \label{fig:sem2sprint1}
\end{figure}

\section{Semester 2 Gantt Chart - Week Level}

\begin{figure}[H]
    \centering
    \includegraphics[width=1\textwidth]{images/sem2_gantt_chart.png}
    \caption{Semester 2 Gantt Chart - Week Level}
    \label{fig:sem2ganttchart}
\end{figure}

\section{Author Document}

\begin{table}[htbp]
\centering
\renewcommand{\arraystretch}{1.5}
\begin{tabular}{>{\raggedright\arraybackslash}p{0.75\textwidth} >{\raggedright\arraybackslash}p{0.18\textwidth}}
\toprule
\rowcolor{lightgray}
\textbf{Task} & \textbf{Assignee} \\
\midrule
Set up GitHub repository under which our decoder solution lives & Christopher \\
\rowcolor{lightgray}
Set up virtual environment for decoder solution & Christopher \\
Create Docker image and container for decoder & Christopher \\
\rowcolor{lightgray}
Run test Docker containing virtualizations & Christopher \\
Study optimizer LION from Google's paper & Tzu Chen \\
\rowcolor{lightgray}
Installing and learning ARC & Tzu Chen \\
Create surface code in Stim to simulate surface code error dataset & Christopher \\
\rowcolor{lightgray}
Generate / simulate Stim datasets & Mara \\
Implement analog readout wrapper & Mara \\
\rowcolor{lightgray}
Connect to the ARC & Arjun \\
Construct embedding scheme (with vectorizing training data and tokenization) & Arjun \\
\rowcolor{lightgray}
Implement proposed transformer error correction architecture with pipeline & Christopher \\
Train on Stim data on ARC & Arjun \\
\rowcolor{lightgray}
Test and evaluate performance on Stim data & Mara \\
Iteratively adjust parameters and model architecture to improve performance on Stim data & Tzu Chen \\
\rowcolor{lightgray}
Implement MWPM baseline error correction on simulated Stim data & Arjun \\
Evaluate MWPM baseline performance on Stim data & Tzu Chen \\
\bottomrule
\end{tabular}
\caption{Project Task Assignments}
\label{tab:tasks}
\end{table}

All group members agree to this breakdown. Chart on next page.

\section{Project Repository}

Our code repository can be found at: \href{https://github.com/PrismaticFlame/qec_decoder_transformer}{QEC Decoder Transformer}

\section{90 Second Teaser Video}

Submitted previously.



\end{document}